% ============================================================
%  Chapter 4 -- Exercises 4.1-4.26
% ============================================================
\rsection{Exercises}

\begin{roadmap}[How to use these]
Chapter 4's exercises divide into a uniform-continuity core, a
counterexample gallery, and a toolbox of functions used throughout the rest
of the book.

\medskip
\raggedright
\textbf{The chain that matters.} \emph{4.9 $\to$ 4.11 $\to$ 4.13} develops
what uniform continuity is really for: it preserves Cauchy sequences, and
therefore lets a function be extended from a dense subset to the whole space
--- the mechanism by which $e^x$, defined on $\Q$ in Exercise 1.6, becomes a
function on $\R$ in Chapter 8.

\smallskip
\textbf{The toolbox.} \emph{4.20--4.22} build the distance function
$\rho_E$ and use it to separate closed sets (normality) --- machinery reused
in Chapters 7 and 10. \emph{4.23--4.24} are convex functions, promised in
the box at 2.17.

\smallskip
\textbf{The gallery.} \emph{4.1, 4.7, 4.16, 4.18, 4.21(b), 4.25(b)} are
counterexamples to remember; 4.18 (Thomae's function) and 4.7
(continuity along every line without continuity) are famous.

\smallskip
\textbf{Alternative proofs.} \emph{4.10} reproves uniform continuity on
compact sets by sequences --- compare it with the cover proof of 4.19 to see
the two faces of compactness side by side.
\end{roadmap}

\begin{enumerate}[label=\textbf{4.\arabic*.}, leftmargin=3.0em, itemsep=9pt,
                  topsep=8pt, labelsep=0.5em]

\item Suppose $f$ is a real function defined on $\R^1$ which satisfies
\[ \lim_{h\to0}\,[f(x+h)-f(x-h)] = 0 \]
for every $x \in \R^1$. Does this imply that $f$ is continuous?\kn{No. The
condition compares the two sides with \emph{each other}, never with $f(x)$
--- so any function with matching one-sided limits but the wrong value
passes: $f = 0$ except $f(0) = 1$. A symmetric condition cannot see an
asymmetric failure. First lesson of the set: check what a limit condition
actually quantifies over.}

\item If $f$ is a continuous mapping of a metric space $X$ into a metric
space $Y$, prove that
\[ f(\overline{E}) \subset \overline{f(E)} \]
for every set $E \subset X$. ($\overline{E}$ denotes the closure of $E$.)
Show, by an example, that $f(\overline{E})$ can be a proper subset of
$\overline{f(E)}$.\kn{Pull back instead of pushing forward:
$\overline{f(E)}$ is closed, so $f^{-1}(\overline{f(E)})$ is closed (corollary
to 4.8) and contains $E$, hence contains $\overline E$. For proper
inclusion, let the closure point escape to infinity: $f(x) = 1/(1+x^2)$ on
$E = \N$ gives $f(\overline E) = f(E)$ but $0 \in \overline{f(E)}$. This is
the exercise promised in Chapter 2's box on why $f^{-1}$ outranks $f$.}

\item Let $f$ be a continuous real function on a metric space $X$. Let
$Z(f)$ (the \emph{zero set} of $f$) be the set of all $p \in X$ at which
$f(p) = 0$. Prove that $Z(f)$ is closed.\kn{One line with the right tool:
$Z(f) = f^{-1}(\{0\})$, and $\{0\}$ is closed. The 4.8 corollary is doing
what it was built for --- and 4.22 below will prove the converse, that every
closed set is some function's zero set.}

\item Let $f$ and $g$ be continuous mappings of a metric space $X$ into a
metric space $Y$, and let $E$ be a dense subset of $X$. Prove that $f(E)$ is
dense in $f(X)$. If $g(p) = f(p)$ for all $p \in E$, prove that $g(p) = f(p)$
for all $p \in X$. (In other words, a continuous function is determined by
its values on a dense subset of its domain.)\kn{The rigidity principle.
For the second part, take $p \in X$, approach it from $E$ (density plus
3.2(d)), and let continuity carry the agreement to the limit. It is why
knowing a continuous function on $\Q$ determines it on $\R$ --- the fact
4.13 will exploit for extensions, and Chapter 8 for $e^x$.}

\item If $f$ is a real continuous function defined on a closed set
$E \subset \R^1$, prove that there exist continuous real functions $g$ on
$\R^1$ such that $g(x) = f(x)$ for all $x \in E$. (Such functions $g$ are
called \emph{continuous extensions} of $f$ from $E$ to $\R^1$.) Show that
the result becomes false if the word ``closed'' is omitted. Extend the
result to vector-valued functions. \emph{Hint:} Let the graph of $g$ be a
straight line on each of the segments which constitute the complement of $E$
(compare Exercise 2.29). The result remains true if $\R^1$ is replaced by
any metric space, but the proof is not so simple.\kn{Exercise 2.29 is the structural input: the open
complement of $E$ is a countable union of disjoint segments, and bridging
each with a chord leaves continuity intact. Without closedness the failure
is immediate --- $1/x$ on $(0,1)$ has nowhere to send $0$. The
vector-valued case is componentwise (4.10).}

\item If $f$ is defined on $E$, the \emph{graph} of $f$ is the set of points
$(x,f(x))$, for $x \in E$. In particular, if $E$ is a set of real numbers
and $f$ is real-valued, the graph of $f$ is a subset of the plane.

Suppose $E$ is compact, and prove that $f$ is continuous on $E$ if and only
if its graph is compact.\kd{A characterisation worth carrying: for compact
domains, continuity of the function equals compactness of the graph. The
forward direction is 4.14 applied to $x \mapsto (x,f(x))$ (continuous by
4.10). The converse is where the work is --- from a sequence $x_n \to x$,
compactness of the graph forces a convergent subsequence of
$(x_n,f(x_n))$ whose limit must be $(x,f(x))$, and 3.20-style bookkeeping
upgrades subsequences to the sequence. Graphs recur at 4.7's pathologies and
in Chapter 9's implicit function theorem.}

\item If $E \subset X$ and if $f$ is a function defined on $X$, the
\emph{restriction} of $f$ to $E$ is the function $g$ whose domain of
definition is $E$, such that $g(p) = f(p)$ for $p \in E$. Define $f$ and $g$
on $\R^2$ by: $f(0,0) = g(0,0) = 0$, $f(x,y) = xy^2/(x^2+y^4)$,
$g(x,y) = xy^2/(x^2+y^6)$ if $(x,y) \ne (0,0)$. Prove that $f$ is bounded on
$\R^2$, that $g$ is unbounded in every neighborhood of $(0,0)$, and that $f$
is not continuous at $(0,0)$; nevertheless, the restrictions of both $f$ and
$g$ to every straight line in $\R^2$ are continuous!\kw{The classic warning
for Chapter 9: continuity along every line through a point does not give
continuity at the point. Approach $f$ along the parabola $x = y^2$ and it
sits at $\tfrac12$; along lines it dies to $0$. Directions are not enough in
dimension two --- curves can sneak in. Remember this when partial
derivatives appear (9.20): their existence is a line-by-line statement, and
this is why it proves so little.}

\item Let $f$ be a real uniformly continuous function on the bounded set
$E$ in $\R^1$. Prove that $f$ is bounded on $E$. Show that the conclusion is
false if the boundedness of $E$ is omitted from the
hypothesis.\kn{Uniform continuity buys finitely many $\delta$-steps across a
bounded set: cover $\overline E$ (compact!) by finitely many $\delta$-balls,
or argue directly --- each step changes $f$ by at most $\varepsilon$, and
boundedness of $E$ caps the number of steps. Compare 4.20(a): plain
continuity could not do this. For the failure, $f(x) = x$ on $\R$.}

\item Show that the requirement in the definition of uniform continuity can
be rephrased as follows, in terms of diameters of sets: To every
$\varepsilon > 0$ there exists a $\delta > 0$ such that
$\operatorname{diam} f(E) < \varepsilon$ for all $E \subset X$ with
$\operatorname{diam} E < \delta$.\kn{Uniform continuity as a statement about
\emph{sets}, with the points quantified away --- 3.9's diameter doing for
functions what it did for Cauchy sequences. This is the form the extension
theorem 4.13 runs on: small sets have small images, uniformly.}

\item Complete the details of the following alternative proof of Theorem
4.19: If $f$ is not uniformly continuous, then for some $\varepsilon > 0$
there are sequences $(p_n)$, $(q_n)$ in $X$ such that
$d_X(p_n,q_n) \to 0$ but $d_Y(f(p_n),f(q_n)) > \varepsilon$. Use Theorem
2.37 to obtain a contradiction.\kd{Worth writing out in full, as the mirror
image of 4.19's cover proof. Negate the uniform statement (one bad
$\varepsilon$, witnesses for every $\delta_n = 1/n$ --- the same
contrapositive harvest as 4.2), compactify: some subsequence
$p_{n_k} \to p$, drag $q_{n_k}$ along since the distance dies, and
continuity at the single point $p$ contradicts the standing gap
$\varepsilon$. Covers versus sequences is the same fork as 2.32 versus
2.41(c) --- Rudin chose covers for the text, Tao's compactness definition
would make this the primary proof.}

\item Suppose $f$ is a uniformly continuous mapping of a metric space $X$
into a metric space $Y$, and prove that $(f(x_n))$ is a Cauchy sequence in
$Y$ for every Cauchy sequence $(x_n)$ in $X$. Use this result to give an
alternative proof of the theorem stated in Exercise 4.13.\kn{The heart of
the uniform/pointwise distinction: a Cauchy sequence's terms crowd each
other at \emph{unpredictable locations}, so only a location-independent
$\delta$ can convert crowding in $X$ to crowding in $Y$. Plain continuity
fails the task: $f(x) = 1/x$ on $(0,1)$ maps the Cauchy sequence $(1/n)$ to
the divergent $(n)$.}

\item A uniformly continuous function of a uniformly continuous function is
uniformly continuous. State this more precisely and prove it.\kn{The relay
of 4.7 again, with both contracts now point-free: $\varepsilon$ buys $\eta$
from $g$ uniformly, $\eta$ buys $\delta$ from $f$ uniformly. Nothing in the
composition argument ever used the base point --- which is the observation
this exercise wants made.}

\item Let $E$ be a dense subset of a metric space $X$, and let $f$ be a
uniformly continuous real function defined on $E$. Prove that $f$ has a
continuous extension from $E$ to $X$ (see Exercise 4.5 for terminology).
(Uniqueness follows from Exercise 4.4.) \emph{Hint:} For each $p \in X$ and
each positive integer $n$, let $V_n(p)$ be the set of all $q \in E$ with
$d(p,q) < 1/n$. Use Exercise 4.9 to show that the intersection of the
closures of the sets $f(V_1(p)),\,f(V_2(p)),\,\dots$ consists of a single
point, say $g(p)$, of $\R^1$. Prove that the function $g$ so defined on $X$
is the desired extension of $f$.

Could the range space $\R^1$ be replaced by $\R^k$? By any compact metric
space? By any complete metric space? By any metric
space?\kd{The most consequential exercise of the chapter. The hint is 3.10(b)
recycled: shrinking sets with shrinking diameters (by 4.9!) whose closures
nest --- their intersection is the value the extension \emph{must} take.
The closing questions locate the true hypothesis: everything works in any
\emph{complete} range space (the nested-closed-sets argument is Exercise
3.21), and completeness is exactly what is needed --- the limit must have
somewhere to live. This machine builds $e^x$ from its rational values
(Chapter 8) and, in functional analysis, extends densely-defined operators.}

\item Let $I = [0,1]$ be the closed unit interval. Suppose $f$ is a
continuous mapping of $I$ into $I$. Prove that $f(x) = x$ for at least one
$x \in I$.\kn{The one-dimensional fixed point theorem, and the standard
first application of 4.23: $g(x) = f(x)-x$ has $g(0) \ge 0$, $g(1) \le 0$,
so the intermediate value theorem plants a zero. One line --- but only
because both endpoints are trapped in $I$. Brouwer's theorem is this
statement in $\R^k$, and needs far heavier tools.}

\item Call a mapping of $X$ into $Y$ \emph{open} if $f(V)$ is an open set in
$Y$ whenever $V$ is an open set in $X$. Prove that every continuous open
mapping of $\R^1$ into $\R^1$ is monotonic.\kn{Contrapositive: a continuous
non-monotonic function has an interior extremum on some $[a,b]$ (attained by
4.16), and then $f((a,b))$ contains its own maximum --- a closed endpoint
--- so is not open. Openness of \emph{images} turns out to be a strong,
rigidity-inducing property; compare the box at 2.2 on why it is not the
definition of continuity.}

\item Let $[x]$ denote the largest integer contained in $x$, that is, $[x]$
is the integer such that $x-1 < [x] \le x$, and let $(x) = x-[x]$ denote the
fractional part of $x$. What discontinuities do the functions $[x]$ and
$(x)$ have?\kn{Simple jumps at every integer, of size $1$, both functions
continuous from the right --- the floor function is 4.31's jump construction
in its simplest form, with $E = \Z$ and all $c_n = 1$. It returns as the
integrator of choice in Chapter 6 (6.16), where its jumps turn integrals
into sums.}

\item Let $f$ be a real function defined on $(a,b)$. Prove that the set of
points at which $f$ has a simple discontinuity is at most countable.
\emph{Hint:} Let $E$ be the set on which $f(x-) < f(x+)$. With each point
$x$ of $E$, associate a triple $(p,q,r)$ of rational numbers such that
\begin{enumerate}[label=(\alph*), leftmargin=2.2em, itemsep=1pt, topsep=3pt]
  \item $f(x-) < p < f(x+)$,
  \item $a < q < t < x$ implies $f(t) < p$,
  \item $x < t < r < b$ implies $f(t) > p$.
\end{enumerate}
The set of all such triples is countable. Show that each triple is
associated with at most one point of $E$. Deal similarly with the other
possible types of simple discontinuities.\kn{4.30's tagging trick without
monotonicity: one rational is no longer enough (jumps can overlap), so each
jump is tagged with a \emph{triple} --- its level plus a rational buffer on
each side --- and 2.13 counts triples. The general principle: to prove a set
countable, inject it into $\Q^n$; the injectivity is where the analysis
lives.}

\item Every rational $x$ can be written in the form $x = m/n$, where
$n > 0$, and $m$ and $n$ are integers without any common divisors. When
$x = 0$, we take $n = 1$. Consider the function $f$ defined on $\R^1$ by
\[ f(x) = \begin{cases} 0 & (x \text{ irrational}),\\[2pt]
   \dfrac1n & \left(x = \dfrac mn\right). \end{cases} \]
Prove that $f$ is continuous at every irrational point, and that $f$ has a
simple discontinuity at every rational point.\kd{Thomae's function ---
continuous exactly on the irrationals. The key: near any point, rationals
with \emph{small} denominators are scarce, so $f(t) < \varepsilon$ once $t$
avoids the finitely many fractions with $n \le 1/\varepsilon$ in a bounded
window. The asymmetry with 4.27(b) is instructive, and one can ask for the
reverse: a function continuous exactly on $\Q$. None exists --- the
continuity set of any function is a countable intersection of open sets,
and $\Q$ is not one, by Baire (Exercise 3.22). That is a satisfying loop:
Baire's theorem answering a question about a single function.}

\item Suppose $f$ is a real function with domain $\R^1$ which has the
intermediate value property: If $f(a) < c < f(b)$, then $f(x) = c$ for some
$x$ between $a$ and $b$. Suppose also, for every rational $r$, that the set
of all $x$ with $f(x) = r$ is closed. Prove that $f$ is continuous.

\emph{Hint:} If $x_n \to x_0$ but $f(x_n) > r > f(x_0)$ for some $r$ and
all $n$, then $f(t_n) = r$ for some $t_n$ between $x_0$ and $x_n$; thus
$t_n \to x_0$. Find a contradiction. (N.\,J.\ Fine, \emph{Amer.\ Math.\
Monthly}, vol.\ 73, 1966, p.\ 782.)\kn{The repaired converse of 4.23: the
intermediate value property alone fails (4.24, 4.27(d)), but adding closed
rational level sets restores continuity. In the hint, the level $t_n$'s
accumulate at $x_0$, so closedness forces $f(x_0) = r$ --- contradiction.
Note how the hypothesis is used only through sequences: the corollary to
4.8 in its sequential clothing.}

\item If $E$ is a nonempty subset of a metric space $X$, define the
distance from $x \in X$ to $E$ by
\[ \rho_E(x) = \inf_{z\in E} d(x,z). \]
\begin{enumerate}[label=(\alph*), leftmargin=2.2em, itemsep=2pt, topsep=3pt]
  \item Prove that $\rho_E(x) = 0$ if and only if $x \in \overline{E}$.
  \item Prove that $\rho_E$ is a uniformly continuous function on $X$, by
        showing that
        \[ |\rho_E(x) - \rho_E(y)| \le d(x,y) \]
        for all $x \in X$, $y \in X$.
\end{enumerate}
\emph{Hint:} $\rho_E(x) \le d(x,z) \le d(x,y) + d(y,z)$, so that
\[ \rho_E(x) \le d(x,y) + \rho_E(y). \]\kn{The workhorse function of metric
topology. Part (a) says $\rho_E$ sees exactly the closure; part (b) says it
is $1$-Lipschitz --- the strongest, cleanest form of uniform continuity,
with $\delta = \varepsilon$. The hint's move (bound by an over-estimate,
then take the inf) is how one always handles infima of distances.}

\item Suppose $K$ and $F$ are disjoint sets in a metric space $X$, $K$ is
compact, $F$ is closed. Prove that there exists $\delta > 0$ such that
$d(p,q) > \delta$ if $p \in K$, $q \in F$. \emph{Hint:} $\rho_F$ is a
continuous positive function on $K$.

Show that the conclusion may fail for two disjoint closed sets if neither is
compact.\kn{4.16 and 4.20 cooperating: $\rho_F$ is positive on $K$ (by
4.20(a), since $F$ is closed and disjoint), and a positive continuous
function on a \emph{compact} set attains a positive minimum. For the
failure take $F = \Z$ and $K = \{n + 1/n\}$: disjoint, closed, distance
$0$. Two closed sets can approach each other at infinity; a compact set has
no infinity to escape to.}

\item Let $A$ and $B$ be disjoint nonempty closed sets in a metric space
$X$, and define
\[ f(p) = \frac{\rho_A(p)}{\rho_A(p)+\rho_B(p)} \qquad (p \in X). \]
Show that $f$ is a continuous function on $X$ whose range lies in $[0,1]$,
that $f(p) = 0$ precisely on $A$ and $f(p) = 1$ precisely on $B$. This
establishes a converse of Exercise 4.3: \emph{Every closed set $A \subset X$
is $Z(f)$ for some continuous real $f$ on $X$.} Setting
\[ V = f^{-1}\left(\left[0,\tfrac12\right)\right), \qquad
   W = f^{-1}\left(\left(\tfrac12,1\right]\right), \]
show that $V$ and $W$ are open and disjoint, and that $A \subset V$,
$B \subset W$. (Thus pairs of disjoint closed sets in metric spaces can be
covered by pairs of disjoint open sets. This property of metric spaces is
called \emph{normality}.)\kd{Urysohn's lemma for metric spaces, built from
Exercise 4.20's bricks in one display. The denominator never vanishes ---
$\rho_A$ and $\rho_B$ cannot both be zero, by 4.20(a) and disjointness ---
and the two preimages under $f$ separate $A$ from $B$ openly, with 4.8
certifying openness. In general topology normality is hard-won; metrics
give it for free, which is one reason this book never needs spaces beyond
metric ones.}

\item A real-valued function $f$ defined in $(a,b)$ is said to be
\emph{convex} if
\[ f(\lambda x + (1-\lambda)y) \le \lambda f(x) + (1-\lambda)f(y) \]
whenever $a < x < b$, $a < y < b$, $0 < \lambda < 1$. Prove that every
convex function is continuous. Prove that every increasing convex function
of a convex function is convex. (For example, if $f$ is convex, so is
$e^f$.)

If $f$ is convex in $(a,b)$ and if $a < s < t < u < b$, show that
\[ \frac{f(t)-f(s)}{t-s} \le \frac{f(u)-f(s)}{u-s} \le
   \frac{f(u)-f(t)}{u-t}. \]\kn{The slope inequality at the end is the real
content --- chords steepen from left to right --- and continuity falls out
of it: near any interior point the difference quotients are wedged between
two fixed chords' slopes, so $f$ is locally Lipschitz. Note ``interior'' is
essential; on $[a,b]$ a convex function may jump at the endpoints. These
inequalities return for the differentiable case in Chapter 5's exercises.}

\item Assume that $f$ is a continuous real function defined in $(a,b)$ such
that
\[ f\left(\frac{x+y}{2}\right) \le \frac{f(x)+f(y)}{2} \]
for all $x,y \in (a,b)$. Prove that $f$ is convex.\kn{Midpoint convexity
plus continuity gives full convexity: iterate midpoints to reach all dyadic
$\lambda = k/2^n$, then let continuity fill in the rest --- density doing
its usual work (Exercise 4.4's principle applied to $\lambda$). Without
continuity the statement fails, but the counterexamples need a Hamel basis,
far beyond this book.}

\item If $A \subset \R^k$ and $B \subset \R^k$, define $A+B$ to be the set
of all sums $\mathbf{x}+\mathbf{y}$ with $\mathbf{x} \in A$,
$\mathbf{y} \in B$.
\begin{enumerate}[label=(\alph*), leftmargin=2.2em, itemsep=3pt, topsep=3pt]
  \item If $K$ is compact and $C$ is closed in $\R^k$, prove that $K+C$ is
        closed.

        \emph{Hint:} Take $\mathbf{z} \notin K+C$, put
        $F = \mathbf{z}-C$, the set of all $\mathbf{z}-\mathbf{y}$ with
        $\mathbf{y} \in C$. Then $K$ and $F$ are disjoint. Choose $\delta$
        as in Exercise 4.21. Show that the open ball with center
        $\mathbf{z}$ and radius $\delta$ does not intersect $K+C$.
  \item Let $\alpha$ be an irrational real number. Let $C_1$ be the set of
        all integers, let $C_2$ be the set of all $n\alpha$ with
        $n \in C_1$. Show that $C_1$ and $C_2$ are closed subsets of
        $\R^1$ whose sum $C_1+C_2$ is not closed, by showing that
        $C_1+C_2$ is a countable dense subset of $\R^1$.\kd{Part (b) is
        the famous fact that $\{m+n\alpha\}$ is dense for irrational
        $\alpha$ --- Weyl's equidistribution in embryo. The proof is
        pigeonhole: the fractional parts $(n\alpha)$ are infinitely many
        points in $[0,1)$, two must fall within $\varepsilon$, and their
        difference generates an $\varepsilon$-mesh. Against (a), it shows
        compactness of one summand is not a removable convenience:
        closed $+$ closed can be dense-and-not-closed.}
\end{enumerate}

\item Suppose $X$, $Y$, $Z$ are metric spaces, and $Y$ is compact. Let $f$
map $X$ into $Y$, let $g$ be a continuous one-to-one mapping of $Y$ into
$Z$, and put $h(x) = g(f(x))$ for $x \in X$. Prove that $f$ is uniformly
continuous if $h$ is uniformly continuous. \emph{Hint:} $g^{-1}$ has compact
domain $g(Y)$, and $f(x) = g^{-1}(h(x))$.

Prove also that $f$ is continuous if $h$ is continuous. Show (by modifying
Example 4.21, or by finding a different example) that the compactness of
$Y$ cannot be omitted from the hypotheses, even when $X$ and $Z$ are
compact.\kn{A cancellation theorem: composing with an injective continuous
$g$ can be undone --- \emph{because} 4.17 makes $g^{-1}$ continuous, indeed
uniformly so by 4.19, and 4.12 composes the uniformities. Every hypothesis
of this chapter's compactness section is consumed in one exercise; the
counterexample without compact $Y$ is 4.21's wrap-around map playing the
role of $g$.}

\end{enumerate}
